\documentclass[11pt]{article}

\usepackage[margin=1in]{geometry}
\usepackage{amsmath,amssymb}
\usepackage{enumitem}
\usepackage{hyperref}
\usepackage{xcolor}
\usepackage{booktabs}
\usepackage{listings}

\hypersetup{colorlinks=true, linkcolor=blue, urlcolor=blue, citecolor=blue}

\lstset{
  basicstyle=\ttfamily\small,
  breaklines=true,
  columns=fullflexible,
}

\title{TOAE209/TOAH209 -- Design and Analysis of Algorithms\\[4pt]
\large Week 10: Practical Exercises}
\author{Foreign Trade University}
\date{}

\begin{document}
\maketitle

\section*{Instructions}

For each problem below there is a starter file
\texttt{exercises/starter\_code\_problemNN.py} containing function signatures with
\texttt{raise NotImplementedError} bodies, and a small \texttt{\_\_main\_\_} block with
tests. Implement the missing functions so the tests pass. A reference solution is
provided in \texttt{solutions/solution\_problemNN.py} -- try the problem yourself before
looking!

All problems use only the Python 3.10+ standard library. Shared helper code lives in
\texttt{starter\_code.py} at the week root: the \texttt{FlowNetwork} class (with
\texttt{add\_edge}, an explicit residual-edge adjacency via \texttt{FlowEdge} objects and
paired reverse edges, a reference \texttt{edmonds\_karp}, and \texttt{copy}/\texttt{value}
utilities), plus the random generators \texttt{generate\_flow\_network} and
\texttt{generate\_bipartite\_graph}.

The 12 problems are grouped into three themes:
\begin{itemize}
    \item \textbf{Flow Networks \& Cuts} (Problems 1--3): the defining conditions of a
    flow, cut capacity, and constructing the residual graph.
    \item \textbf{Max-Flow Algorithms} (Problems 4--8): BFS augmenting paths,
    Edmonds-Karp, DFS Ford-Fulkerson, and extracting a minimum cut from a maximum flow.
    \item \textbf{Verification \& Applications} (Problems 9--12): bipartite matching,
    multiple sources/sinks, edge-disjoint paths (Menger), and integrality -- the
    reductions and structural guarantees that make flow so widely applicable.
\end{itemize}

\newpage
\section{Flow Networks \& Cuts}

\subsection*{Problem 1 -- Flow Feasibility Checker}

Implement \texttt{verify\_flow(network, flow)}. Given a \texttt{FlowNetwork} and a
candidate \texttt{flow} (a dict mapping each original edge \texttt{(u, v)} to a flow
amount), check that the assignment is a valid flow: (1) every edge satisfies $0 \le f(e)
\le c(e)$ (capacity constraints), and (2) flow is conserved at every node other than the
source and sink (flow in $=$ flow out). Return \texttt{(True, value)} when valid, where
\texttt{value} is the net flow leaving the source; return \texttt{(False, 0.0)} otherwise.
A flow that references an edge not present in the network is invalid.

\textbf{Example.} On the network with edges $s\!\to\!a(3)$, $s\!\to\!b(2)$, $a\!\to\!b(1)$,
$a\!\to\!t(2)$, $b\!\to\!t(3)$, the assignment $\{sa{:}2, sb{:}2, at{:}2, bt{:}2\}$ is
valid with value $4$.

\subsection*{Problem 2 -- Cut Capacity}

Implement \texttt{cut\_capacity(network, S)}. Given an $s$-$t$ cut described by a set
\texttt{S} (which must contain the source and not the sink), return the cut capacity: the
sum of capacities of edges going \emph{from} a node in \texttt{S} \emph{to} a node not in
\texttt{S}. Edges from $T$ back into $S$ do not contribute. Raise \texttt{ValueError} if
\texttt{S} omits the source or contains the sink.

\textbf{Example.} On the network above, \texttt{cut\_capacity(G, \{"s"\})} $= 5$ (edges
$s\!\to\!a$ and $s\!\to\!b$).

\subsection*{Problem 3 -- Build the Residual Graph for a Given Flow}

Implement \texttt{build\_residual\_graph(network, flow)}. For each original edge $(u,v)$
with capacity $c$ and flow $f$: include a forward residual edge $u\!\to\!v$ with residual
capacity $c-f$ when $c-f>0$, and a backward residual edge $v\!\to\!u$ with residual
capacity $f$ when $f>0$. Return an adjacency dict mapping each node to a \emph{sorted} list
of \texttt{(neighbour, residual\_capacity)} pairs, omitting zero-capacity residual edges.

\textbf{Hint.} The backward edges are exactly what let an augmenting-path algorithm cancel
previously routed flow.

\newpage
\section{Max-Flow Algorithms}

\subsection*{Problem 4 -- BFS for an Augmenting Path}

Implement \texttt{bfs\_augmenting\_path(network)}. Using the current edge flows stored in
the \texttt{FlowNetwork} (each \texttt{FlowEdge} exposes \texttt{residual\_capacity()}),
run a breadth-first search from the source and return a \emph{shortest} (fewest-edges)
augmenting path as a list of nodes \texttt{[s, ..., t]}, or \texttt{None} if the sink is
unreachable in the residual graph. If the source equals the sink, return \texttt{[s]}.

\subsection*{Problem 5 -- Edmonds-Karp Maximum Flow}

Implement \texttt{edmonds\_karp(network)}. Repeatedly find a shortest augmenting path by
BFS; if none exists, stop. Otherwise compute the bottleneck (minimum residual capacity
along the path) and push that much flow along it, updating forward and backward residual
edges (use \texttt{FlowEdge.push}). Mutate the edge flows in place and return the max-flow
value. This is Ford-Fulkerson with the BFS rule, giving the $O(VE^2)$ bound.

\textbf{Example.} \texttt{edmonds\_karp(G) == 5.0} on the running example; on the
unit-capacity ``diamond'' ($s\!\to\!a,s\!\to\!b,a\!\to\!b,a\!\to\!t,b\!\to\!t$) it returns
\texttt{2.0}, demonstrating flow cancellation via the residual back-edges.

\subsection*{Problem 6 -- Ford-Fulkerson with DFS Augmenting Paths}

Implement \texttt{ford\_fulkerson(network)} using \emph{depth-first} search to find
augmenting paths instead of BFS. Any augmenting path yields a correct max flow with integer
capacities (DFS does not give the $O(VE^2)$ bound, but terminates on these small integer
instances). Mutate the edge flows in place and return the value. It should agree with
Edmonds-Karp on every instance.

\subsection*{Problem 7 -- Verify Edmonds-Karp and Ford-Fulkerson Agree}

Implement \texttt{verify\_algorithms\_agree(num\_trials, n, seed)}: for each of
\texttt{num\_trials} random networks (\texttt{generate\_flow\_network(n, seed=seed +
trial)}), compute the max flow with both your Edmonds-Karp (Problem 5) and your DFS
Ford-Fulkerson (Problem 6), and return \texttt{True} iff the values match on every trial.
Because the max-flow value is \emph{unique} (it equals the min-cut), any two correct
augmenting-path methods must agree. Use \texttt{network.copy()} so each algorithm runs on
its own copy (the flow mutations are in place).

\subsection*{Problem 8 -- Extract the Minimum Cut from a Max Flow}

Implement \texttt{min\_cut(network)}: compute a maximum flow, then let \texttt{S} be the
set of nodes reachable from the source in the \emph{final} residual graph (BFS along
positive-residual edges). Return \texttt{(S, max\_flow\_value)}; the cut $(S, V\setminus
S)$ is a minimum cut whose capacity equals the max-flow value (Max-Flow Min-Cut). Also
implement \texttt{verify\_max\_flow\_min\_cut(num\_trials, n, seed)} returning \texttt{True}
iff, on every random network, \texttt{cut\_capacity(G, S)} equals the max-flow value and
the sink is not in \texttt{S}.

\newpage
\section{Verification \& Applications}

\subsection*{Problem 9 -- Bipartite Matching via Max Flow}

Implement \texttt{max\_bipartite\_matching(left, right, edges)} by reducing to max flow:
add a super-source with a unit-capacity edge to each left vertex, a unit-capacity edge for
each allowed pair (left $u$ to right $v$), and a unit-capacity edge from each right vertex
to a super-sink. The max-flow value equals the maximum matching size; by integrality the
flow selects a valid set of matching edges. Label left vertices \texttt{("L", u)} and right
vertices \texttt{("R", v)} so the two sides stay distinct.

\textbf{Example.} With \texttt{left=right=3} and edges $\{(0,0),(1,1),(2,2)\}$ the matching
size is $3$; with all of left pointing only at right vertex $0$, it is $1$.

\subsection*{Problem 10 -- Multiple Sources/Sinks via Super-Source/Super-Sink}

Implement \texttt{multi\_source\_sink\_max\_flow(nodes, edges, sources, sinks)} where each
edge is a tuple \texttt{(u, v, capacity)}. Introduce a super-source connected to every
original source by an infinite-capacity edge, and a super-sink fed by every original sink
via an infinite-capacity edge; the max flow between them equals the maximum total flow the
sources can deliver to the sinks. Return that value.

\textbf{Example.} With sources $\{0,1\}$, sink $\{3\}$, and edges $0\!\to\!2(3),
1\!\to\!2(2), 2\!\to\!3(10)$, the answer is $5$; lowering $2\!\to\!3$ to capacity $4$ makes
it $4$.

\subsection*{Problem 11 -- Edge-Disjoint Paths via Unit-Capacity Max Flow (Menger)}

Implement \texttt{max\_edge\_disjoint\_paths(nodes, directed\_edges, s, t)}: the maximum
number of pairwise edge-disjoint $s$-$t$ paths. By Menger's theorem (the edge form), this
equals the max $s$-$t$ flow when \emph{every} edge has capacity $1$; with integral flow,
each unit traces one edge-disjoint path. Return the count.

\textbf{Example.} On $s\!\to\!a\!\to\!t$ and $s\!\to\!b\!\to\!t$ the answer is $2$; on a
single chain $s\!\to\!a\!\to\!b\!\to\!t$ it is $1$.

\subsection*{Problem 12 -- Integrality of the Computed Flow}

Implement \texttt{flow\_is\_integral(network)}: assuming a flow has been computed, return
\texttt{True} iff every original forward edge carries an integer flow. Then implement
\texttt{verify\_integrality(num\_trials, n, seed)}: for each random integer-capacity
network, compute the max flow and check the flow is integral on every edge and the total
value is itself an integer. Return \texttt{True} iff this holds on every trial.

\textbf{Background.} The Integrality Theorem guarantees this: starting from the all-zero
flow with integer capacities, every augmentation pushes an integer bottleneck, so the flow
stays integral throughout -- which is exactly what makes the $0/1$ matching and
edge-disjoint-path reductions (Problems 9 and 11) work.

\newpage
\section{LeetCode Practice}

After completing the problems above, practise this week's flow / min-cut / matching
\emph{modelling} skills on the following \href{https://leetcode.com}{LeetCode} problems.
Note that \emph{pure} max-flow problems are uncommon on LeetCode; the problems below are
ones that are naturally modelled as flow, min-cut, or bipartite matching (and several are
solved in practice by matching DP or BFS + min-cut intuition rather than a full
Ford-Fulkerson implementation). Difficulty is LeetCode's own label.

\begin{itemize}
  \item \href{https://leetcode.com/problems/maximum-number-of-achievable-transfer-requests/}{Maximum Bipartite Matching $\to$ Maximum Number of Achievable Transfer Requests} -- \emph{Hard} -- flow / enumeration.
  \item \href{https://leetcode.com/problems/maximum-students-taking-exam/}{Maximum Students Taking Exam} -- \emph{Hard} -- matching / bitmask.
  \item \href{https://leetcode.com/problems/minimum-cost-to-connect-two-groups-of-points/}{Minimum Cost to Connect Two Groups of Points} -- \emph{Hard} -- assignment / min-cost matching.
  \item \href{https://leetcode.com/problems/campus-bikes-ii/}{Campus Bikes II} -- \emph{Medium} -- assignment (min-cost matching).
  \item \href{https://leetcode.com/problems/pizza-with-3n-slices/}{Pizza With 3n Slices} -- \emph{Hard} -- matching-style DP.
  \item \href{https://leetcode.com/problems/escape-the-spreading-fire/}{Escape the Spreading Fire} -- \emph{Hard} -- BFS + min-cut intuition.
\end{itemize}

\end{document}
